\documentclass[12pt,a4paper]{article}
\usepackage{amsmath,amssymb,amsthm}
\usepackage{hyperref}
\usepackage{geometry}
\geometry{margin=2.5cm}
\usepackage{enumitem}
\usepackage{booktabs}

\newtheorem{theorem}{Theorem}[section]
\newtheorem{lemma}[theorem]{Lemma}
\newtheorem{corollary}[theorem]{Corollary}
\newtheorem{proposition}[theorem]{Proposition}
\theoremstyle{definition}
\newtheorem{definition}[theorem]{Definition}
\theoremstyle{remark}
\newtheorem{remark}[theorem]{Remark}

\title{Gravitational Wave Speed Equals Light Speed:\\
A Structural Proof from the Common Ledger Substrate}

\author{Jonathan Washburn\\
Recognition Science Research Institute\\
Austin, Texas\\
\texttt{washburn.jonathan@gmail.com}}

\date{March 2026}

\begin{document}
\maketitle

\begin{abstract}
The multi-messenger observation of the binary neutron star merger
GW170817 constrains the gravitational-wave propagation speed
$c_{\mathrm{grav}}$ to be equal to the speed of light $c$ to within
$|c_{\mathrm{grav}} - c|/c < 10^{-15}$.  This raises a structural
question: is this near-equality approximate or exact?  We prove that
in Recognition Science (RS), $c_{\mathrm{grav}} = c$ \emph{exactly},
as a consequence of both sectors propagating on the same discrete
ledger substrate with the same tick rate $\tau_0$.  The speed of any
signal on the ledger is bounded above by one voxel per tick:
$v_{\max} = \ell_0/\tau_0 = c$.  Since both electromagnetic and
gravitational perturbations propagate as recognition events on this
common substrate, their propagation speeds are identically $c$.  The
speed ratio $c_{\mathrm{grav}}/c = 1$ is not an empirical coincidence
but a structural identity.  \end{abstract}

%% ============================================================
\section{Introduction}
%% ============================================================

General relativity predicts that gravitational waves propagate at the
speed of light.  This was spectacularly confirmed on 17 August 2017,
when LIGO/Virgo detected the gravitational-wave signal GW170817 and
the Fermi Gamma-ray Burst Monitor detected the associated short
gamma-ray burst GRB~170817A, arriving $1.7\;\mathrm{s}$ apart after
traveling $\sim 40\;\mathrm{Mpc}$~\cite{GW170817}.  The constraint:
\begin{equation}
  \frac{|c_{\mathrm{grav}} - c|}{c} < 5 \times 10^{-16}.
\end{equation}

While this is an extraordinary observational confirmation, it leaves
open the foundational question: \emph{why} does gravity propagate at
exactly $c$?  In GR, this follows from the structure of the field
equations, but the field equations themselves are postulated.

Recognition Science~\cite{RS_Axioms} derives the exact equality
$c_{\mathrm{grav}} = c$ from a deeper structural principle: both
sectors propagate on the same ledger substrate.

%% ============================================================
\section{The Ledger Speed Limit}
\label{sec:limit}
%% ============================================================

\begin{definition}[Ledger Speed Limit]
\label{def:c}
In the discrete ledger with voxel size $\ell_0$ and tick period
$\tau_0$, the causal propagation speed is
\begin{equation}
  c = \frac{\ell_0}{\tau_0}.
\end{equation}
A recognition event at voxel $(i, j, k)$ at tick $t$ can influence
its nearest neighbor at tick $t+1$; no influence can travel faster
than one voxel per tick.  This maximum speed defines $c$.
\end{definition}

\begin{definition}[Gravitational Perturbation Speed]
\label{def:cgrav}
The \emph{gravitational propagation speed} $c_{\mathrm{grav}}$ is
the speed at which a linearized perturbation $h_{\mu\nu}$ of the
emergent metric propagates through the ledger, as measured by
counting the number of ticks required for the perturbation to travel
one voxel.
\end{definition}

\begin{remark}
Definition~\ref{def:cgrav} is \emph{not} the same as
Definition~\ref{def:c} by construction; $c_{\mathrm{grav}}$ is a
derived quantity that must be computed from the RS dynamics.  The
equality $c_{\mathrm{grav}} = c$ is a non-trivial result, not a
definitional tautology.
\end{remark}

%% ============================================================
\section{The Structural Proof}
\label{sec:proof}
%% ============================================================

\begin{theorem}[Common-Substrate Speed Equality]
\label{thm:forced}
Both electromagnetic and gravitational perturbations are massless
recognition events on the same discrete ledger.  Consequently:
\begin{equation}
  c_{\mathrm{grav}} = c.
\end{equation}
\end{theorem}

\begin{proof}
The RS ledger is a single discrete structure with one causal cone:
a perturbation can propagate at most one voxel per tick, giving the
upper bound $c = \ell_0/\tau_0$ for any signal.

Electromagnetic waves propagate as oscillating recognition events
in the $U(1)$ ledger channel.  In the massless limit, the photon
saturates the causal bound: $c_{\mathrm{EM}} = c$.

Gravitational waves propagate as linearized perturbations
$h_{\mu\nu}$ of the emergent spacetime metric.  In RS, the metric
emerges from the $J$-cost field on the ledger (the Einstein field
equations arise from stationarity of the RS action; see the
companion paper on EFE emergence).  The linearized gravitational
wave equation in RS-native units takes the form
$\Box h_{\mu\nu} = 0$, where $\Box = \partial_t^2 - \nabla^2$
is the d'Alembertian with $c = 1$.  Massless solutions of
$\Box h_{\mu\nu} = 0$ propagate at speed $c$.

Both wave equations ($\Box A_\mu = 0$ and $\Box h_{\mu\nu} = 0$)
live on the \emph{same} lattice with the \emph{same} tick rate and
the \emph{same} $c$.  Therefore $c_{\mathrm{grav}} = c_{\mathrm{EM}} = c$.
\end{proof}

\begin{theorem}[Speed Ratio Unity]
\label{thm:ratio}
$c_{\mathrm{grav}}/c = 1$ exactly.
\end{theorem}

\begin{proof}
Immediate from Theorem~\ref{thm:forced}: $c_{\mathrm{grav}} = c$,
so $c_{\mathrm{grav}}/c = 1$.
\end{proof}

\begin{corollary}[Graviton Mass]
\label{cor:graviton}
The graviton mass is exactly zero in RS.  A massive graviton would
propagate at $c_{\mathrm{grav}} < c$ (subluminal dispersion), violating
the ledger speed identity.  Since the ledger enforces
$c_{\mathrm{grav}} = c$ structurally, no graviton mass term is
permitted.  This is consistent with the experimental bound
$m_g < 1.76 \times 10^{-23}\;\mathrm{eV}/c^2$
from GW170817~\cite{deRham2017}.
\end{corollary}

%% ============================================================
\section{Comparison with GR and Alternatives}
\label{sec:compare}
%% ============================================================

\begin{center}
\renewcommand{\arraystretch}{1.3}
\begin{tabular}{@{}lcc@{}}
\toprule
\textbf{Framework} & $c_{\mathrm{grav}} = c$? &
\textbf{Status} \\
\midrule
General Relativity & Yes (from field equations) & Postulated \\
Massive gravity & No ($c_{\mathrm{grav}} < c$) & Falsified
by GW170817 \\
Brans--Dicke / scalar-tensor & Depends on parameters & Constrained \\
Recognition Science & Yes (from common substrate) & Derived \\
\bottomrule
\end{tabular}
\end{center}

The RS derivation is more fundamental than GR's: in GR, the equality
follows from the postulated Lorentz symmetry of the field equations.
In RS, it follows from the \emph{physical fact} that there is one
lattice, one tick rate, and hence one speed limit.

%% ============================================================
\section{Observational Consistency}
\label{sec:obs}
%% ============================================================

The GW170817 observation constrains
$|c_{\mathrm{grav}} - c|/c < 5 \times 10^{-16}$~\cite{GW170817}.
RS predicts \emph{exact} equality: $c_{\mathrm{grav}}/c = 1$ with no
correction at any order.  This is consistent with, and stronger than,
the observational constraint.

Future multi-messenger events with longer baselines (cosmological
distances) can tighten this bound further.  RS predicts that no
deviation will ever be found, regardless of precision.

%% ============================================================
\section{Conclusion}
\label{sec:conclusion}
%% ============================================================

Gravitational waves propagate at exactly the speed of light because
both gravity and electromagnetism are recognition events on the same
discrete ledger.  The common substrate enforces a common speed limit
$c = \ell_0/\tau_0$.  The speed ratio $c_{\mathrm{grav}}/c = 1$ is
a structural identity, not an empirical near-coincidence.

%% ============================================================
\begin{thebibliography}{99}

\bibitem{GW170817}
B.~P.~Abbott \emph{et al.} (LIGO/Virgo and Fermi GBM
Collaborations), ``Gravitational Waves and Gamma-Rays from a Binary
Neutron Star Merger: GW170817 and GRB~170817A,'' Astrophys.\ J.\
Lett.\ \textbf{848}, L13 (2017).

\bibitem{deRham2017}
C.~de~Rham, J.~T.~Deskins, A.~J.~Tolley, and S.-Y.~Zhou,
``Graviton Mass Bounds,'' Rev.\ Mod.\ Phys.\ \textbf{89}, 025004
(2017).

\bibitem{RS_Axioms}
J.~Washburn, ``Recognition Science: Derivation of the Cost
Functional and Forcing Chain,'' Recognition Science Research
Institute Technical Report (2025).

\end{thebibliography}

\end{document}
